\documentclass[10pt,a4paper]{amsart}
\usepackage[margin=19mm]{geometry}
\usepackage{amsmath,amssymb,mathtools}
\usepackage[scr=rsfs,cal=euler]{mathalfa}
\usepackage{xfrac}
\usepackage[unicode,hidelinks,bookmarks=false]{hyperref}
\newcommand{\ie}{i.e.\ }
\newcommand{\cf}{cf.\ }
\newcommand{\resp}{respectively\ }
\newcommand{\Subsection}{\subsection}
\providecommand{\nobreakdash}{\nobreak\hbox{-}}
\setlength{\emergencystretch}{2em}
\multlinegap0pt
\usepackage{bm}
\usepackage{bbm}
\usepackage{dsfont}
\usepackage{xspace}
\usepackage{cancel}
\usepackage{mathtools}
\usepackage{amsthm}
\makeatletter
\@ifundefined{prop}{\newtheorem{prop}{Prop}}{}
\makeatother
\makeatletter
\newcommand{\K}{\mathbb K}
\expandafter\ifx\csname algoproof\endcsname\relax
\newenvironment{algoproof}{\vspace{1em}\noindent\textit{Algorithmic Proof.}}{\hfill\qedsymbol\vspace{1em}}
\fi
\makeatother
\title[A linear-time algorithm for Chow decompositions]{A linear-time algorithm for Chow decompositions}
\author{Alexander Taveira Blomenhofer, Benjamin Lovitz}
\date{Source: arXiv:2509.10450v1 (CC BY 4.0)}
\begin{document}
\maketitle

\noindent\emph{Source and licence.} A linear-time algorithm for Chow decompositions. Version arXiv:2509.10450v1 (CC BY 4.0), distributed under the Creative Commons Attribution 4.0 International licence, as recorded in the arXiv licence metadata for this exact version and shown on the abstract page https://arxiv.org/abs/2509.10450v1. Full rights notice: licenses/arxiv-2509.10450-CC-BY-4.0.txt.\par\smallskip

\noindent\textit{Editorial scope.} Counted unit: the Prop whose source label is \texttt{prop:grouping-test} in the article (source0.tex lines 539--552, source label \texttt{prop:grouping-test}), delivered together with the article's own complete original proof of it (source0.tex lines 553--592, one \texttt{proof} environment of 40 lines). No mathematics is written, repaired or re-derived: statement and proof are the article's own text line for line. Required context retained uncounted: none. Every definition, display and label of those lines is retained unchanged. Omitted from this package and not redistributed: 1--538, 593--823. Nothing inside a retained excerpt is omitted or altered: every retained body is delivered exactly as the article's lines are written, so no omission receipt of any kind accompanies this package. Citations: the retained excerpts carry no citation command at all, so no bibliography entry of the article is transcribed and no reference is redistributed. Reference adaptation: none was needed and none was made; every reference inside the delivered unit resolves inside it, so no unresolved reference or citation is produced. Printed slips kept literal: no printed word, symbol, hypothesis or number of the counted statement or of its proof was altered. Layout and class: the article's own class is \texttt{amsart} with packages including inputenc, amssymb, amscd, thmtools, todonotes, hyperref, enumitem, float, cleveref, stmaryrd, pifont; the wrapper is the lane's pinned \texttt{amsart} editorial wrapper at 10pt on A4, which restates the theorem-like environments the retained text needs and adds the article's own macro definitions that the retained text needs (\texttt{\textbackslash K}), with extra wrapper packages (bm, bbm, dsfont, xspace, cancel, mathtools). The delivered numbering is the unit's own. Assets: no retained span contains a figure environment, a graphics-inclusion command, a PDF-page inclusion, a table or a TikZ picture; no image file is copied into this package, the package declares no asset dependency and no image file is redistributed with it. Licence: version arXiv:2509.10450v1 of this article is distributed under the Creative Commons Attribution 4.0 International licence, as recorded in the arXiv licence metadata for this exact version and shown on the abstract page https://arxiv.org/abs/2509.10450v1. A full rights notice is delivered at licenses/arxiv-2509.10450-CC-BY-4.0.txt. Characters: every retained excerpt is pure ASCII, so the delivered engine, XeLaTeX with its default Latin Modern text font, reports no missing character.
\begin{prop}\label{prop:grouping-test}
	Let $T=e_1 \cdots e_{2d+1}$, let $\partial=e_1+\dots + e_{2d}$ and $ \partial_{d+1}=e_1+\dots + e_{2d+1} $,  let $y=\prod_{j=1}^d (e_{2j}-e_{2j-1})$, let $\mu=e_1+\dots+e_{2d-1}+e_{2d+1}$, and let $z=(e_{2d+1}-e_{2d})\prod_{j=1}^{d-1} (e_{2j}-e_{2j-1})$. Then
	\begin{align}\label{eq:evector_check}
		T_{\partial} y =0,
	\end{align}
	and similarly,
	\begin{align}\label{eq:evector_check2}
		T_{\mu} z =0.
	\end{align}
	Furthermore,
	\begin{align}\label{eq:contraction_nonzero}
		z^T T_{\partial_{2d+1}} y \neq 0.
	\end{align}
\end{prop}

\begin{proof}
	Denote by $ \sigma(R) $ the number of odd elements in a set $ R $. Note that
	\begin{align*}
		y=\sum_{\substack{R \subseteq [2d],\: |R|=d\\ \forall  j\in [d]\colon |R \cap \{2j-1,2j\}|=1}} (-1)^{\sigma(R)} e_R.
	\end{align*}
	From here, it is easily calculated that for the contraction of $ T $ by $ y $, we have
	\begin{align*}
		(I_n^{\otimes d+1} \otimes y^{T})T= \alpha \; e_{2d+1} y
	\end{align*}
	for some non-zero scalar $\alpha \in \K^{\times}$. We first verify~\eqref{eq:evector_check}. Note that
	\begin{align*}
		T_{\partial} y = \alpha \sum_{\substack{R \subseteq [2d]\\ |R|=d\\ |R \cap \{2j-1,2j\}|=1 \forall  j\in [d]}} (-1)^{\sigma(R)} w_R,
	\end{align*}
	where
	\begin{align*}
		w_R=e_{2d+1}\sum_{i \in R} e_{R \setminus \{i\}}.
	\end{align*}
	
	%Using the expression~\eqref{eq:Tdel} for $T_{\partial}$, we have
	%% is clear from the proof of Proposition~\ref{prop:geom_mult} that 
	%\begin{align*}
	%T_{\partial} \left(\prod_{j=1}^d (x_{2j}-x_{2j-1})\right)= \sum_{\substack{R \subseteq [2d]\\ |R|=d\\ |R \cap \{2j-1,2j\}|=1 \forall  j\in [d]}} (-1)^{\sum R} w_R,
	%%\sum_{\substack{S \subseteq [2d]\\ |S|=d}} w_S,
	%\end{align*}
	%where
	%\begin{align*}
	%w_R=x_{2d+1}\sum_{i \in [2d]\setminus R} x_{[2d]\setminus (R \cup \{i\})}.
	%\end{align*}
	Let $S \subseteq [2d]$ be a subset of size $d-1$. If $|S \cap \{2j,2j-1\}|=2$ for some $j \in [d]$, then none of the $w_R$ are supported on $S \cup \{2d+1\}$. Otherwise, the only possibility is that $|S \cap \{2j,2j-1\}|=1$ for all but one index $j' \in [d]$, and for this index $|S \cap \{2j',2j'-1\}|=0$. In this case, there are only two choices of $R$ for which $w_{R}$ is supported on $S \cup \{2d+1\}$, namely $R=S \cup \{2j'\}$ and $R=S \cup \{2j'-1\}$. The coefficients of $w_R$ for these two choices are additive inverses, hence $T_{\partial} y$ is not supported on $S \cup \{2d+1\}$. This verifies~\eqref{eq:evector_check}. By symmetry,~\eqref{eq:evector_check2} also holds.
	
	It remains to verify~\eqref{eq:contraction_nonzero}. Note that
	\begin{align*}
	T_{\partial_{2d+1}} y = \alpha y,
	\end{align*}
	so
	\begin{align*}
	z^T T_{\partial_{2d+1}} y=\alpha \; z^T y= \alpha \; 2^{d-1} (e_{2d+1}-e_{2d-1})^T (e_{2d}-e_{2d-1})= \alpha \; 2^{d-1} \neq 0.
	\end{align*}
	This completes the proof.
\end{proof}

\end{document}
